\section*{Supplementary Materials}

\subsection*{Full family of pairwise comparisons}

Table~\ref{tab:supp-holm} lists all 14 planned pairwise comparisons on
which the Holm-Bonferroni correction in Table~\ref{tbl:statistical} was
based. Nine of these comparisons are highlighted in the main text.

\begin{table}[htbp]
\centering
\caption{Full family of pairwise Wilcoxon signed-rank tests with
Holm-Bonferroni correction.}\label{tab:supp-holm}
\resizebox{\linewidth}{!}{\begin{tabular}{@{}lccccccc@{}}
\toprule
Comparison & MAE A & MAE B & HL $\Delta$ & 95\% CI & $r_{rb}$ & $p$ raw & $p$ Holm \\
\midrule
Best RF (auditable ML) vs Best DL & 0.0200 & 0.0268 & -0.0042 & [-0.0164, 0.0003] & -0.49 & 0.0934 & 1.0000 \\
Best RF (spectrogram) vs Best WST & 0.0200 & 0.0225 & -0.0011 & [-0.0039, 0.0018] & -0.26 & 0.3755 & 1.0000 \\
RF dB vs Ridge dB & 0.0210 & 0.0248 & -0.0050 & [-0.0097, -0.0012] & -0.60 & 0.0335 & 0.4360 \\
RF WST vs Ridge WST & 0.0248 & 0.0310 & -0.0088 & [-0.0147, 0.0015] & -0.49 & 0.0934 & 1.0000 \\
RF dB vs Shallow MLP features & 0.0210 & 0.0495 & -0.0249 & [-0.0412, -0.0127] & -1.00 & 3.05e-5 & 4.27e-4 \\
RF Vib-dB vs ResNetVibCNN Vib-dB & 0.0221 & 0.0268 & -0.0032 & [-0.0086, -0.0001] & -0.54 & 0.0577 & 0.6921 \\
RF AE/Vib dB-z vs RF Vib-dB-z & 0.0200 & 0.0223 & -0.0003 & [-0.0022, 0.0008] & -0.06 & 0.8603 & 1.0000 \\
RF dB-z vs RF dB & 0.0200 & 0.0210 & -0.0006 & [-0.0031, 0.0017] & -0.18 & 0.5619 & 1.0000 \\
Best RF vs Best LightGBM & 0.0200 & 0.0250 & -0.0008 & [-0.0115, 0.0009] & -0.26 & 0.3755 & 1.0000 \\
ResNetVibCNN vs BilinearFusion & 0.0268 & 0.0317 & -0.0029 & [-0.0094, 0.0034] & -0.31 & 0.2979 & 1.0000 \\
LightGBM Vib-WST vs ResNetVibCNN Vib-dB & 0.0225 & 0.0268 & -0.0023 & [-0.0157, 0.0006] & -0.41 & 0.1591 & 1.0000 \\
RF AE-dB vs ResNetAECNN & 0.0403 & 0.0300 & 0.0045 & [-0.0016, 0.0266] & 0.34 & 0.2522 & 1.0000 \\
ResNetAECNN vs ChannelAttentionCNN & 0.0300 & 0.0406 & -0.0098 & [-0.0229, 0.0029] & -0.37 & 0.2114 & 1.0000 \\
ResNetVibCNN vs ResNetAECNN & 0.0268 & 0.0300 & -0.0002 & [-0.0123, 0.0093] & -0.01 & 0.9799 & 1.0000 \\
\bottomrule
\end{tabular}
}
\end{table}

\subsection*{Nested grouped sensitivity for tree models}

The canonical benchmark remains post-selection. As a sensitivity analysis,
representative tree-model configurations were re-evaluated with outer 16-fold
LOGO and inner grouped five-fold selection among a small fixed hyperparameter
grid. This is not a nested comparison of every model family, but it tests
whether the observed tree-model ordering is highly dependent on the fixed tree
hyperparameters used in the canonical analysis. The legacy log-mel run and
the subsequent cache-matched reruns are both retained below: their discrepancy
is a provenance finding, not evidence that either representation is uniquely
superior.

\begin{table}[htbp]
\centering
\caption{Nested LOGO sensitivity analysis for representative tree-model
configurations. Each outer test condition is excluded from inner selection.}
\label{tab:supp-nested-logo}
\begin{tabular}{@{}llcc@{}}
\toprule
Model & Configuration & Outer LOGO MAE (\(\upmu\)m) & Fold SD (\(\upmu\)m) \\
\midrule
Random forest & AE-dB + Vib-dB & 0.0213 & 0.0441 \\
Random forest & AE-dB-z + Vib-dB-z (current rerun) & 0.0207 & 0.0347 \\
Random forest & AE-log-mel + Vib-log-mel (legacy) & 0.0455 & 0.0590 \\
Random forest & AE-log-mel + Vib-log-mel (current rerun) & 0.0206 & 0.0480 \\
LightGBM & AE-dB + Vib-dB & 0.0296 & 0.0514 \\
LightGBM & Vib-dB & 0.0281 & 0.0514 \\
Random forest & AE/Vib features + physics + PP & 0.0493 & 0.0613 \\
\bottomrule
\end{tabular}
\end{table}

The legacy log-mel result is not reproduced by the current cache-matched
rerun. The matched
dB-z and log-mel analyses give 0.0207 and 0.0206~\(\upmu\)m, respectively,
under the same outer/inner fold design and RF grid. This confirms the central
limitation of the point-estimate ranking: the leading dB-z and log-mel
representations form an observed group rather than a uniquely selected winner.
The RF grid was \{100 trees/unlimited depth, 200 trees/unlimited depth,
100 trees/depth 10, 200 trees/depth 8\}; every candidate used random seed
42. Five grouped inner folds selected the lowest mean validation MAE within
each outer fold, after which the selected model was refitted on all
outer-training conditions. Per-fold selections are archived in
\path{nested_logo_dbz_matched_folds.csv} and
\path{nested_logo_logmel_matched_folds.csv}.

\begin{longtable}{@{}lrrrrr@{}}
\caption{Per-outer-fold selections for the current nested RF sensitivity.
The five inner folds are grouped by grinding condition. Selection minimises
mean inner-fold MAE; the selected RF is refitted on the same 14 outer-training
conditions used by the fixed RF, while the archived validation condition
remains excluded, and is evaluated once on the held-out outer condition.}
\label{tab:supp-nested-rf-selections}\\
\toprule
Representation & Test condition & Trees & Maximum depth & Inner MAE & Outer MAE \\
\midrule
\endfirsthead
\toprule
Representation & Test condition & Trees & Maximum depth & Inner MAE & Outer MAE \\
\midrule
\endhead
dB-z & 1 & 100 & 10 & 0.0282 & 0.0052 \\
dB-z & 2 & 100 & 10 & 0.0254 & 0.0035 \\
dB-z & 3 & 100 & None & 0.0211 & 0.0097 \\
dB-z & 4 & 200 & None & 0.0203 & 0.0080 \\
dB-z & 5 & 200 & 8 & 0.0226 & 0.0033 \\
dB-z & 6 & 200 & 8 & 0.0189 & 0.0059 \\
dB-z & 7 & 100 & 10 & 0.0170 & 0.1402 \\
dB-z & 8 & 100 & None & 0.0283 & 0.0727 \\
dB-z & 9 & 200 & None & 0.0170 & 0.0090 \\
dB-z & 10 & 100 & None & 0.0202 & 0.0107 \\
dB-z & 11 & 100 & None & 0.0215 & 0.0086 \\
dB-z & 12 & 200 & 8 & 0.0166 & 0.0096 \\
dB-z & 13 & 100 & None & 0.0215 & 0.0126 \\
dB-z & 14 & 100 & 10 & 0.0210 & 0.0170 \\
dB-z & 15 & 200 & 8 & 0.0276 & 0.0098 \\
dB-z & 16 & 200 & 8 & 0.0226 & 0.0052 \\
log-mel & 1 & 200 & 8 & 0.0313 & 0.0051 \\
log-mel & 2 & 100 & 10 & 0.0276 & 0.0051 \\
log-mel & 3 & 200 & None & 0.0216 & 0.0091 \\
log-mel & 4 & 100 & 10 & 0.0211 & 0.0117 \\
log-mel & 5 & 100 & None & 0.0198 & 0.0048 \\
log-mel & 6 & 200 & None & 0.0213 & 0.0084 \\
log-mel & 7 & 200 & None & 0.0213 & 0.2061 \\
log-mel & 8 & 100 & None & 0.0282 & 0.0130 \\
log-mel & 9 & 100 & 10 & 0.0164 & 0.0115 \\
log-mel & 10 & 100 & 10 & 0.0265 & 0.0094 \\
log-mel & 11 & 200 & None & 0.0220 & 0.0052 \\
log-mel & 12 & 100 & 10 & 0.0160 & 0.0108 \\
log-mel & 13 & 100 & 10 & 0.0272 & 0.0062 \\
log-mel & 14 & 100 & 10 & 0.0259 & 0.0057 \\
log-mel & 15 & 200 & None & 0.0299 & 0.0102 \\
log-mel & 16 & 100 & 10 & 0.0207 & 0.0074 \\
\bottomrule
\end{longtable}


\begin{table}[htbp]
\centering
\caption{Canonical whole-condition assignments generated by successive
draws from \texttt{numpy.random.RandomState(42)}. Training comprises all
conditions other than the listed test and validation conditions.}
\label{tab:supp-logo-assignments}
\begin{tabular}{@{}ccc@{\qquad}ccc@{}}
\toprule
Fold/test & Validation & Training count & Fold/test & Validation & Training count \\
\midrule
1 & 8 & 14 & 9 & 7 & 14 \\
2 & 5 & 14 & 10 & 11 & 14 \\
3 & 14 & 14 & 11 & 3 & 14 \\
4 & 16 & 14 & 12 & 7 & 14 \\
5 & 12 & 14 & 13 & 11 & 14 \\
6 & 9 & 14 & 14 & 11 & 14 \\
7 & 14 & 14 & 15 & 8 & 14 \\
8 & 5 & 14 & 16 & 5 & 14 \\
\bottomrule
\end{tabular}
\end{table}


\begin{longtable}{@{}p{2.25cm}p{2.15cm}ccp{2.45cm}p{2.6cm}c@{}}
\caption{RF experiment-provenance map. The canonical split manifest is
\texttt{reports/evidence/tables/canonical\_logo\_assignments.csv}; cache
hashes are given in Table~\ref{tab:supp-cache-provenance}. Unless stated
otherwise, current RFs use 200 trees, depth eight, one minimum sample per
leaf, random state 42, no fold-wise spectral scaler, Python 3.11.11,
NumPy 2.1.3, and scikit-learn 1.6.1.}
\label{tab:supp-rf-provenance}\\
\toprule
Experiment ID & Input/role & Train conditions & Per-sample z & Split or selection & Distinguishing implementation & Mean MAE \\
\midrule
\endfirsthead
\toprule
Experiment ID & Input/role & Train conditions & Per-sample z & Split or selection & Distinguishing implementation & Mean MAE \\
\midrule
\endhead
canonical-db-z & AE/Vib dB-z ranking & 14 & Yes & Canonical seed-42; historical & Exact package environment unavailable & 0.0200 \\
canonical-log-mel & AE/Vib log-mel ranking & 14 & No & Canonical seed-42; historical & Exact package environment unavailable & 0.0201 \\
robustness-base & AE/Vib dB-z ablation baseline & 15 & Yes & Test condition only excluded & Regenerated suite; explains 0.0202 baseline & 0.0202 \\
current-fixed-db-z & AE/Vib dB-z primary current & 14 & Yes & Canonical seed-42 & Validation condition remains excluded & 0.0210 \\
current-fixed-db & AE/Vib dB current & 14 & No & Canonical seed-42 & Validation condition remains excluded & 0.0213 \\
nested-db-z & AE/Vib dB-z sensitivity & 14 & Yes & Five inner grouped folds & Selected model refitted on same 14 conditions & 0.0207 \\
nested-log-mel & AE/Vib log-mel sensitivity & 14 & No & Five inner grouped folds & Selected model refitted on same 14 conditions & 0.0206 \\
rf-seed-repeat-14 & AE/Vib dB-z equal-budget seed sensitivity & 14 & Yes & Canonical test and validation excluded & Seeds 42--46; seed-42 0.0210, five-seed mean 0.0203; used for paired model-family sensitivity & 0.0203 \\
rf-seed-repeat-15 & AE/Vib dB-z deployment-style seed sensitivity & 15 & Yes & Test condition only excluded & Historical seeds 42--46; seed-42 0.0201, five-seed mean 0.0195; excluded from model-family tests & 0.0195 \\
modality-mask-suite & Modality and test-time masking & 15 & Yes & Test condition only excluded & Same regenerated 0.0202 full-input baseline & 0.0202 \\
retrained-band & Retrained band removal & 14 & Yes & Canonical seed-42 & Features zeroed before fit and test & 0.0210 \\
rf-oob-uq & OOB-scaled uncertainty only & 15 & Yes & Test condition only excluded & Fold scaler; 200 trees; depth 8; min leaf 2 & -- \\
rf-oof-shap & OOF attribution & 14 & Yes & Canonical seed-42 & Approximate path-dependent TreeSHAP on held-out passes & 0.0210 \\
latency-rf-db-z & Fold-complete timing & 14 & Yes & All 16 canonical repeat-0 checkpoints & Pooled 16,000-call timing; complete checkpoint hashes in \texttt{edge\_latency\_benchmark\_per\_fold.csv} & 0.0200 \\
\bottomrule
\end{longtable}

The 14-condition seed repeat uses the same training-data budget as the neural
baselines. The retained 15-condition seed repeat and robustness-suite values
do not reproduce the 14-condition current fixed RF because they use the
validation condition for fitting; they are not used for paired model-family
inference.
The latency row pools timings over the same 16 canonical repeat-0 checkpoints
that produce the historical 16-fold benchmark aggregate, so its latency and
accuracy now refer to the same fitted-model family. Modality, masking, and OOB analyses use 15 non-test
conditions and are therefore method-specific diagnostics rather than exact
refits of the recommended 14-condition pipeline.

The historical 14-comparison family reads its 16 fold records per
model--configuration from \texttt{full\_results\_logo\_only.csv} (SHA-256
\nolinkurl{b91bae9e26c8ff3f2f9e965e570d93e6cf6d3c72b103bdda41f6f79463ba0ca0}).
The canonical split manifest has SHA-256
\nolinkurl{5b19e7babe5f7de3f38f064dc4797fc37adf31aaab1d9f7cd6f3f7b425971ea5},
and the generated statistical table has SHA-256
\nolinkurl{f902e18c29105f4a09f7a9b64765eb55b0fe15608f0eb9344093f318ae220cca}.
The 240 fold-level prediction files used across the 15 unique
model--configuration sources in that family are listed with complete hashes
in \texttt{canonical\_prediction\_manifest.csv}; the manifest itself has
SHA-256
\nolinkurl{96cfe7989a2666b6add9023570a690e7e2de0dd42c235818df339c4c1bc5c286}.


\begin{table}[htbp]
\centering
\caption{Deep validation-condition rotation sensitivity. Each model is
retrained with fixed default architecture and training settings
($\mathrm{lr}=10^{-3}$, weight decay $10^{-4}$, maximum 100 epochs,
patience 15); only the whole-condition validation assignment changes.
This is not a re-tuned or nested comparison of the canonical deep models.}
\label{tab:supp-deep-validation-rotation}
\begin{tabular}{@{}llccc@{}}
\toprule
Model & Validation assignment & Mean MAE (\(\upmu\)m) & Fold SD & Mean epochs \\
\midrule
ResNetVibCNN & Seed-42 draw & 0.0393 & 0.0519 & 29.8 \\
ResNetVibCNN & Cyclic next condition & 0.0471 & 0.0675 & 23.8 \\
ResNetAECNN & Seed-42 draw & 0.0440 & 0.0660 & 22.6 \\
ResNetAECNN & Cyclic next condition & 0.0435 & 0.0686 & 23.7 \\
\bottomrule
\end{tabular}
\end{table}

\begin{table}[htbp]
\centering
\caption{Physics-feature variance audit across the 16 LOGO training folds.
Five of the nominal 44 aggregate columns are constant in every training
fold and are mapped to zero by fold-wise standardisation.}
\label{tab:supp-physics-audit}
\begin{tabular}{@{}lcc@{}}
\toprule
Quantity & Count & Detail \\
\midrule
Constructed aggregate columns & 44 & 11 indicators $\times$ 4 aggregates \\
Nonconstant campaign columns & 39 & Retain variation across passes \\
Constant in every LOGO training fold & 5 & burst\_rate\_narrow\_min, burst\_rate\_broad\_min, $ec$\_std, $bid$\_std, $st$\_std \\
\bottomrule
\end{tabular}
\end{table}

\subsection*{Per-sample normalisation ablation}

Table~\ref{tab:supp-psnorm-ablation} compares the recommended random
forest trained on AE-dB-z + Vib-dB-z with and without
per-sample/channel standardisation. Removing the per-sample operation while
leaving the otherwise unscaled RF spectral pipeline unchanged produces only
a small numerical change in overall LOGO MAE.

\begin{table}[htbp]
\centering
\caption{Strict per-sample normalisation ablation for the recommended
random forest. Both variants use 200 trees, maximum depth 8, minimum one
sample per leaf, and no fold-wise spectrogram scaling, matching the
canonical RF pipeline; the ``without'' variant omits only per-sample
standardisation. Both variants use the canonical seed-42 LOGO
validation-condition exclusions; experiment ID:
\texttt{psnorm-strict-logo-seed42-v4}.}\label{tab:supp-psnorm-ablation}
\begin{tabular}{@{}lcccc@{}}
\toprule
Variant & Overall MAE (\(\upmu\)m) & Median condition MAE (\(\upmu\)m) & Max condition MAE (\(\upmu\)m) \\
\midrule
With per-sample norm & 0.0210 & 0.0093 & 0.1478 \\
Without per-sample norm & 0.0213 & 0.0093 & 0.1863 \\
\bottomrule
\end{tabular}
\end{table}

\begin{table}[htbp]
\centering
\caption{Retrained frequency-band ablation for the 14-condition current
benchmark RF. Selected
per-sample dB-z feature bins are set to zero before fitting and testing in
every seed-42 LOGO fold. Intervals are descriptive 95\% bootstrap intervals
for the mean paired condition-level MAE difference versus the full dB-z RF.}
\label{tab:supp-retrained-band-ablation}
\begin{tabular}{@{}lccc@{}}
\toprule
Variant & Mean MAE (\(\upmu\)m) & Difference vs full (\(\upmu\)m) & 95\% CI \\
\midrule
Full dB-z & 0.0210 & 0.00000 & [0.00000, 0.00000] \\
Without AE >1~MHz & 0.0210 & 0.00001 & [-0.00075, 0.00082] \\
Without Vib 2--15~kHz & 0.0208 & -0.00021 & [-0.00085, 0.00034] \\
Without both bands & 0.0212 & 0.00024 & [-0.00142, 0.00222] \\
Without Vib >15~kHz & 0.0285 & 0.00755 & [0.00149, 0.01532] \\
Without Vib 16--24~kHz & 0.0234 & 0.00239 & [-0.00318, 0.00846] \\
\bottomrule
\end{tabular}
\end{table}

\begin{table}[htbp]
\centering
\caption{Strict pre-normalisation vibration-band ablation for the
14-condition current RF. Bands are deleted from the dB cache, each
sample/channel z-score is recomputed using retained frequencies only, and the
RF is refitted. Intervals are descriptive paired condition-bootstrap 95\%
intervals versus the recomputed full dB-z baseline.}
\label{tab:supp-prez-band-ablation}
\begin{tabular}{@{}lccc@{}}
\toprule
Variant & Mean MAE (\(\upmu\)m) & Difference vs full (\(\upmu\)m) & 95\% CI \\
\midrule
Full dB-z recomputed from dB & 0.0209 & 0.00000 & [0.00000, 0.00000] \\
Remove Vib >15~kHz before z-score & 0.0255 & 0.00455 & [0.00063, 0.01066] \\
Remove Vib 16--24~kHz before z-score & 0.0254 & 0.00447 & [-0.00035, 0.01050] \\
\bottomrule
\end{tabular}
\end{table}

\begin{table}[htbp]
\centering
\caption{Fresh-process runtime-memory diagnostic for the latency candidates.
Cached inputs are loaded before the baseline RSS measurement; the reported
increment is measured after model loading and one warm-up inference. It is
not a peak-memory or end-to-end deployment measurement.}
\label{tab:supp-runtime-memory}
\begin{tabular}{@{}llc@{}}
\toprule
Model & Configuration & Resident delta (MB) \\
\midrule
Random forest & AE-dB-z + Vib-dB-z & 5.90 \\
Random forest & AE-dB + Vib-dB & 5.27 \\
LightGBM & Vib-dB-z & 7.24 \\
ResNetVibCNN & Vib-dB & 15.39 \\
ResNetAECNN & AE-dB & 13.83 \\
Bilinear fusion & AE-dB + Vib-dB + physics + PP & 20.23 \\
\bottomrule
\end{tabular}
\end{table}

\subsection*{Pass-position sensitivity}

The canonical cache averages all local dB spectrogram maps within a pass.
For a temporal-position sensitivity analysis, every pass was partitioned
into nonempty early, middle, and late thirds of its own local-map axis;
the two modality maps in each third were averaged and per-plane z-scored
before fitting the same 200-tree, depth-eight RF under the 16 seed-42
LOGO folds. The ordered multi-window row concatenates all three
third-level descriptors before fitting; it is not a mean of independently
trained predictions. This is an archived-pass analysis, not an online
fixed-window or fixed-distance experiment.

\begin{table}[htbp]
\centering
\caption{Pass-position sensitivity for the recommended RF using
AE-dB-z + Vib-dB-z inputs. Values are mean $\pm$ standard deviation of
the 16 condition-level LOGO MAEs.}
\label{tab:supp-pass-position}
\begin{tabular}{@{}lcc@{}}
\toprule
Local-map segment & Mean MAE (\(\upmu\)m) & Fold SD (\(\upmu\)m) \\
\midrule
Early third & 0.0213 & 0.0351 \\
Middle third & 0.0204 & 0.0349 \\
Late third & 0.0198 & 0.0335 \\
Ordered early + middle + late & 0.0195 & 0.0338 \\
\bottomrule
\end{tabular}
\end{table}

\subsection*{Run order and measured roughness}

The campaign nominally contained 320 passes. Condition~1, sample~1 was
excluded because the corresponding AE/vibration intermediate files were
corrupt or missing, leaving 319 valid passes. Figure~\ref{fig:supp-runorder}
shows the measured surface roughness in experimental run order; Condition~7
is highlighted. A run-order effect cannot be ruled out because the 16
conditions were run in a single randomized sequence and wheel-wear or
thermal-drift trends may be confounded with the condition effect.

\begin{figure}[htbp]
\centering
\includegraphics[width=\linewidth]{supp_run_order.png}
\caption{Run-order plot of measured $R_a$ for the nominal 320 passes;
one pass was excluded only from model evaluation because its
AE/vibration files were missing. The red \textsf{x} marks
Condition~7.}\label{fig:supp-runorder}
\end{figure}

\begin{figure}[htbp]
\centering
\includegraphics[width=\linewidth]{supp_runorder_diagnostics.png}
\caption{Run-order diagnostics for the best random forest (AE-dB-z +
Vib-dB-z). \textbf{a}, Measured $R_a$. \textbf{b}, RF prediction residual
($\hat y - y$). \textbf{c}, AE broadband RMS. \textbf{d}, Tri-axial vibration RMS.
Condition~7 is highlighted in red; it does not coincide with a clear
run-order trend in sensor RMS.}\label{fig:supp-runorder-diagnostics}
\end{figure}

\begin{figure}[htbp]
\centering
\includegraphics[width=\linewidth]{supp_full_residual_diagnostics.png}
\caption{Residual diagnostics for the recommended random forest
(AE-dB-z + Vib-dB-z). Residuals are plotted against \textbf{a}, measured
\(R_a\); \textbf{b}, predicted \(R_a\); \textbf{c}, wheel speed;
\textbf{d}, workpiece speed; \textbf{e}, depth of cut; \textbf{f}, run order;
\textbf{g}, AE RMS; \textbf{h}, vibration RMS; and \textbf{i}, distance
outside the training-target support. Condition~7 samples are highlighted in red.}
\label{fig:supp-full-residual-diagnostics}
\end{figure}

Across the same samples, RF residuals are most strongly associated with
measured \(R_a\) (Spearman \(\rho=-0.484\) for signed residual) and
training-target support distance (\(\rho=-0.355\) for signed residual and
\(\rho=0.429\) for absolute residual). Absolute residuals also increase
with predicted \(R_a\) (\(\rho=0.421\)). Correlations with process
settings and sensor RMS are weaker: absolute-residual Spearman
correlations are 0.159 for wheel speed, 0.083 for workpiece speed,
\(-0.155\) for depth of cut, 0.220 for run order, \(-0.162\) for AE RMS,
and \(-0.124\) for vibration RMS. Samples outside the training-target
support have mean residual \(-0.1228~\upmu\)m and mean absolute residual
0.1271~\(\upmu\)m, compared with 0.0057 and 0.0117~\(\upmu\)m for
in-support samples, respectively.

% Flush the preceding supplementary floats before starting the next table block.
\clearpage

\begin{table}[htbp]
\centering
\caption{Cluster-aware diagnostics for full LOGO MC-dropout uncertainty.
Pass-level quantities are descriptive because the 319 passes are clustered
within 16 conditions. Cluster-bootstrap intervals resample conditions.}
\label{tab:supp-clustered-uncertainty}
\begin{tabular}{@{}lccc@{}}
\toprule
Metric & Pass-level & Cluster-bootstrap 95\% CI & Condition means ($n=16$) \\
\midrule
Pearson $r$ & 0.488 & [0.202, 0.951] & 0.492 \\
Spearman $\rho$ & 0.594 & [0.252, 0.817] & 0.732 \\
AUROC (top 20\% error) & 0.818 & [0.471, 0.992] & --- \\
\bottomrule
\end{tabular}
\end{table}

\subsection*{Per-condition LOGO MAE for top models}

Figure~\ref{fig:supp-per-condition-mae} shows the per-condition LOGO MAE
for the best random forest, the best spectrogram CNN
(ResNetVibCNN on Vib-dB), and the best gradient-boosted model (LightGBM on
Vib-dB-z). The historical fold-level Vib-WST artifact is unavailable and is
therefore not drawn. Condition~7 is a large error spike for all three models,
indicating a robust prediction failure rather than a model-specific
failure.

\begin{figure}[htbp]
\centering
\includegraphics[width=\linewidth]{supp_per_condition_mae.png}
\caption{Per-condition LOGO MAE for three top models. The vertical dashed
line marks Condition~7.}\label{fig:supp-per-condition-mae}
\end{figure}

\subsection*{Physics-informed feature definitions}

The 44-dimensional physics-informed feature vector is obtained by
aggregating the indicators listed in Table~\ref{tab:physics_features}
over time (mean, standard deviation, maximum, minimum). The indicators
are computed from the raw AE and vibration signals as follows.

\begin{table}[htbp]
\centering
\caption{Physics-informed feature indicators.}\label{tab:physics_features}
\begin{tabular}{@{}p{2.6cm}p{7.2cm}ll@{}}
\toprule
Feature & Formula / definition & Units & Agg. \\
\midrule
wavelet\_energy\_narrow & $\sum_{f\in[50,400]\,\text{kHz}} |W_f|^2$ (complex Morlet CWT, \texttt{cmor3-3}, scales 1--99) & a.u. & mean, std, max, min \\
wavelet\_energy\_broad & $\sum_{f\in[150,2000]\,\text{kHz}} |W_f|^2$ capped at the AE Nyquist frequency of 2~MHz (complex Morlet CWT, \texttt{cmor3-3}, scales 1--99) & a.u. & mean, std, max, min \\
burst\_rate\_narrow & Count of peaks above $4\sigma$ in the 50--400~kHz AE band, minimum spacing $0.1$~ms, divided by window length & counts/s & mean, std, max, min \\
burst\_rate\_broad & Count of peaks above $4\sigma$ in the 150--2000~kHz AE band capped at 2~MHz Nyquist, minimum spacing $0.1$~ms, divided by window length & counts/s & mean, std, max, min \\
$EC$ & $F_t=1.2 a_p^{0.8} v_s^{-0.2}$; $e_c=\frac{F_t v_s}{a_p v_w 214}$ and $EC=e_c v_{w,\text{ref}}$, where $v_{w,\text{ref}}$ linearly maps $v_w\in[20,80]$ to $[0.85,1.15]$; the cached feature is $EC$ & a.u. & mean, std, max, min \\
$bid$ & $a_p / 30$ (normalised depth-of-cut indicator) & a.u. & mean, std, max, min \\
$st$ & $\frac{T_{\max}}{800}$, from the implemented transient semi-infinite heating ODE in Eqs.~\eqref{eq:ec-implemented}--\eqref{eq:st-implemented} & a.u. & mean, std, max, min \\
env\_kurtosis\_x & $\text{kurtosis}(|x(t)+j\hat x(t)|)$ (Hilbert envelope) & a.u. & mean, std, max, min \\
env\_kurtosis\_y & $\text{kurtosis}(|y(t)+j\hat y(t)|)$ & a.u. & mean, std, max, min \\
env\_kurtosis\_z & $\text{kurtosis}(|z(t)+j\hat z(t)|)$ & a.u. & mean, std, max, min \\
$mag$ & $\sqrt{x^2(t)+y^2(t)+z^2(t)}$ & a.u. & mean, std, max, min \\
\bottomrule
\end{tabular}
\end{table}

Here $v_w$ (r/min), $v_s$ (m/s), and $a_p$ (\(\upmu\)m) are the workpiece
speed, wheel speed, and grinding depth recorded in the process-parameter
file. The pre-processing used for the cached burst-rate features is
encoded in cached intermediate files. The archive does not preserve the
complete filter design and edge-handling settings, so these cached burst
values can be reused but not independently reconstructed from raw signals.

The temperature indicator $st$ used in the cached ML features is the
following implemented heuristic. First compute
\begin{equation}
\label{eq:ec-implemented}
F_t=1.2a_p^{0.8}v_s^{-0.2},\qquad
e_c=\frac{F_t v_s}{a_p v_w\,214},\qquad
EC=e_c\left[0.85+0.30\frac{v_w-20}{60}\right].
\end{equation}
Here \(e_c\) is an intermediate unadjusted proxy. The cached
44-dimensional vector stores the speed-adjusted quantity \(EC\), not both
\(e_c\) and \(EC\); \(EC\) is also the quantity used by the temperature
heuristic below. The constant 214 is the workpiece outer diameter in mm, retained in the
cached empirical formula. Define $q=0.8v_s EC$, thermal conductivity
$K=46\times10^{-3}$ W/(mm K), and volumetric heat capacity
$\rho c_p=3.8\times10^{-3}$ J/(mm$^3$ K). The implementation solves
\begin{equation}
\label{eq:st-implemented}
\frac{dT}{dt}=\frac{q\sqrt{a_p}}{K\sqrt{\pi(K/(\rho c_p))t}},\qquad
T(10^{-6}\,\mathrm{s})=25^\circ\mathrm{C},\qquad
T_{\max}=T(10\,\mathrm{s}),\qquad st=\frac{T_{\max}}{800},
\end{equation}
with SciPy's Radau integrator. The documented reference calculation gives
\(T_{\max}\approx26.9^\circ\mathrm{C}\) and \(st\approx0.0336\).
This is a cached process-derived severity heuristic, not a calibrated
grinding-temperature measurement or a dimensionally exact thermal model.

\subsection*{Design balance}

Table~\ref{tab:supp-balance} shows the pairwise occurrence counts of the
factor levels in the 16-run design. Wheel speed ($v_s$) is fully crossed
with workpiece speed ($v_w$) and depth of cut ($a_p$), but $v_w$ and
$a_p$ are not fully crossed: each level pair appears once or twice.

\begin{table}[htbp]
\centering
\caption{Pairwise factor-level occurrence counts in the 16-run
fractional factorial design.}\label{tab:supp-balance}
\begin{tabular}{@{}lcccc@{}}
\toprule
$v_s$ (m/s) & $v_w=20$ & $v_w=40$ & $v_w=60$ & $v_w=80$ \\
\midrule
25 & 1 & 1 & 1 & 1 \\
30 & 1 & 1 & 1 & 1 \\
35 & 1 & 1 & 1 & 1 \\
40 & 1 & 1 & 1 & 1 \\
\bottomrule
\end{tabular}
\vspace{0.5em}

\begin{tabular}{@{}lcccc@{}}
\toprule
$v_s$ (m/s) & $a_p=20$ & $a_p=30$ & $a_p=40$ & $a_p=50$ \\
\midrule
25 & 1 & 1 & 1 & 1 \\
30 & 1 & 1 & 1 & 1 \\
35 & 1 & 1 & 1 & 1 \\
40 & 1 & 1 & 1 & 1 \\
\bottomrule
\end{tabular}
\vspace{0.5em}

\begin{tabular}{@{}lcccc@{}}
\toprule
$v_w$ (r/min) & $a_p=20$ & $a_p=30$ & $a_p=40$ & $a_p=50$ \\
\midrule
20 & 1 & 2 & 1 & 0 \\
40 & 0 & 1 & 2 & 1 \\
60 & 1 & 0 & 1 & 2 \\
80 & 2 & 1 & 0 & 1 \\
\bottomrule
\end{tabular}
\end{table}

\subsection*{MC-dropout uncertainty subset (earlier exploratory case study)}

The 59-sample MC-dropout case study corresponds to all passes from
Conditions 1, 2, and 6 selected by random seed 42, as summarized in Table~\ref{tab:supp-mcdropout-subset}. It was an exploratory
pilot used before the full 16-fold LOGO evaluation and is retained here
only for comparison; the current main result is the full-LOGO evaluation
in Table~\ref{tab:supp-mcdropout-full-logo}. Condition~7 is not
represented in this subset.

\begin{table}[htbp]
\centering
\caption{Composition of the 59-sample MC-dropout uncertainty subset.}
\label{tab:supp-mcdropout-subset}
\begin{tabular}{@{}lcc@{}}
\toprule
Condition & Passes in subset & Passes available \\
\midrule
1 & 19 & 19 \\
2 & 20 & 20 \\
6 & 20 & 20 \\
\midrule
Total & 59 & 59 \\
\bottomrule
\end{tabular}
\end{table}

\subsection*{MC-dropout full LOGO evaluation}

Table~\ref{tab:supp-mcdropout-full-logo} reports per-condition results
for the full 16-fold LOGO MC-dropout evaluation of ResNetVibCNN on
Vib-dB. Aggregate coverage is 50.2\% across 319 held-out
samples; Condition~7 has 0\% coverage.

\begin{table}[htbp]
\centering
\caption{Full 16-fold LOGO MC-dropout coverage per condition
(ResNetVibCNN on Vib-dB).}\label{tab:supp-mcdropout-full-logo}
\begin{tabular}{@{}lcccc@{}}
\toprule
Condition & Passes & Coverage (\%) & Mean width ($\upmu$m) & MAE ($\upmu$m) \\
\midrule
1 & 19 & 0.0\% & 0.05729 & 0.04341 \\
2 & 20 & 80.0\% & 0.02403 & 0.00740 \\
3 & 20 & 95.0\% & 0.03644 & 0.00851 \\
4 & 20 & 15.0\% & 0.19246 & 0.13304 \\
5 & 20 & 80.0\% & 0.02896 & 0.00807 \\
6 & 20 & 85.0\% & 0.03205 & 0.00768 \\
7 & 20 & 0.0\% & 0.04197 & 0.24722 \\
8 & 20 & 10.0\% & 0.05477 & 0.08767 \\
9 & 20 & 20.0\% & 0.09187 & 0.08234 \\
10 & 20 & 20.0\% & 0.02393 & 0.01608 \\
11 & 20 & 70.0\% & 0.03388 & 0.01341 \\
12 & 20 & 55.0\% & 0.01557 & 0.00843 \\
13 & 20 & 70.0\% & 0.02237 & 0.00955 \\
14 & 20 & 85.0\% & 0.01958 & 0.00485 \\
15 & 20 & 40.0\% & 0.06440 & 0.03791 \\
16 & 20 & 75.0\% & 0.02675 & 0.00774 \\
All & 319 & 50.2\% & 0.04787 & 0.04521 \\
\bottomrule
\end{tabular}
\end{table}


\begin{table}[htbp]
\centering
\caption{Fold-specific diagnostic comparing deterministic checkpoint
predictions with the mean of 50 stochastic MC-dropout forwards for
ResNetVibCNN on Vib-dB. Prediction shift is the mean absolute difference
between deterministic and stochastic-mean predictions.}
\label{tab:supp-mcdropout-deterministic}
\begin{tabular}{@{}rrrrr@{}}
\toprule
Condition & Deterministic MAE & Stochastic MAE & Prediction shift & Coverage (\%) \\
\midrule
1 & 0.0389 & 0.0434 & 0.0048 & 0.0 \\
2 & 0.0073 & 0.0074 & 0.0009 & 80.0 \\
3 & 0.0087 & 0.0085 & 0.0011 & 95.0 \\
4 & 0.1312 & 0.1330 & 0.0059 & 15.0 \\
5 & 0.0064 & 0.0081 & 0.0019 & 80.0 \\
6 & 0.0076 & 0.0077 & 0.0008 & 85.0 \\
7 & 0.2489 & 0.2472 & 0.0022 & 0.0 \\
8 & 0.0845 & 0.0877 & 0.0034 & 10.0 \\
9 & 0.0787 & 0.0823 & 0.0041 & 20.0 \\
10 & 0.0155 & 0.0161 & 0.0008 & 20.0 \\
11 & 0.0154 & 0.0134 & 0.0022 & 70.0 \\
12 & 0.0084 & 0.0084 & 0.0004 & 55.0 \\
13 & 0.0093 & 0.0096 & 0.0008 & 70.0 \\
14 & 0.0040 & 0.0048 & 0.0016 & 85.0 \\
15 & 0.0343 & 0.0379 & 0.0043 & 40.0 \\
16 & 0.0090 & 0.0077 & 0.0020 & 75.0 \\
\bottomrule
\end{tabular}
\end{table}


\subsection*{Model--configuration inventory}

The complete inventory of evaluated model--configuration pairs, together
with each pair's mean LOGO MAE and whether it entered the LOGO-only
ranking, is provided in the file
\path{reports/evidence/tables/supp_model_inventory.csv} in the
repository. A concise matrix of model variants versus representation
families is shown in Table~\ref{tab:supp-inventory}.

\begin{table}[htbp]
\centering
\caption{Evaluated model--representation matrix. A checkmark indicates
that at least one model--input configuration in that representation family
was evaluated under the LOGO protocol. It does not imply that every fused
fixed-learner configuration shown in Table~\ref{tbl:controlled-representations}
exists.}
\label{tab:supp-inventory}
\footnotesize
\begin{tabular}{@{}lccccccccc@{}}
\toprule
Model family & PP & Physics & AE features & Vib features & AE-dB & Vib-dB & AE-dB-z & Vib-dB-z & WST \\
\midrule
Ridge & \checkmark & \checkmark & \checkmark & \checkmark & \checkmark & \checkmark & \checkmark & \checkmark & \checkmark \\
Random forest & \checkmark & \checkmark & \checkmark & \checkmark & \checkmark & \checkmark & \checkmark & \checkmark & \checkmark \\
LightGBM & \checkmark & \checkmark & \checkmark & \checkmark & \checkmark & \checkmark & \checkmark & \checkmark & \checkmark \\
MLP & \checkmark & \checkmark & \checkmark & \checkmark & \checkmark & \checkmark & \checkmark & \checkmark & \checkmark \\
ResNetVibCNN & -- & -- & -- & -- & -- & \checkmark & -- & -- & -- \\
ResNetAECNN & -- & -- & -- & -- & \checkmark & -- & -- & -- & -- \\
ChannelAttentionCNN & -- & -- & -- & -- & \checkmark & -- & -- & -- & -- \\
TrajectoryCNN & \checkmark & -- & -- & -- & \checkmark & \checkmark & -- & -- & -- \\
ResNetFusion & -- & -- & -- & -- & \checkmark & \checkmark & \checkmark & \checkmark & -- \\
BilinearFusion & -- & -- & -- & -- & \checkmark & \checkmark & \checkmark & \checkmark & -- \\
\bottomrule
\end{tabular}
\end{table}

\subsection*{Cache provenance and exact representation definitions}

Tables~\ref{tab:supp-cache-definitions} and~\ref{tab:supp-cache-provenance}
list the six principal cached time--frequency representations, their exact
transform chains and shapes, and compact checksum prefixes. Complete SHA-256
digests remain in the public machine-readable manifest.

\begin{table}[htbp]
\centering
\caption{Exact transformation definitions for the principal time--frequency representations. Internal cache identifiers are retained here only for provenance.}
\label{tab:supp-cache-definitions}
\small
\resizebox{\linewidth}{!}{%
\begin{tabular}{@{}l p{2.8cm} p{2.7cm} p{3.0cm} p{2.8cm} l@{}}
\toprule
Cache & Source quantity & Frequency transform & Amplitude transform & Normalization & Shape \\
\midrule
ae\_spec & Arithmetic mean of 2,910 local AE dB maps & Linear STFT bins (NFFT 598, 300 bins) & Magnitude -> dB, clipped at -80 dB, mean over local-map axis & None in cache; no fold-wise scaling in RF spectral runs & (319, 2, 300, 47) \\
vib\_spec & Arithmetic mean of 2,910 local vibration dB maps & Linear STFT bins (NFFT 512, 257 bins) & Magnitude -> dB, clipped at -80 dB, mean over local-map axis & None in cache; no fold-wise scaling in RF spectral runs & (319, 3, 257, 13) \\
ae\_logspec & Pass-mean AE dB cache & None & Per-sample z-score of dB spectrogram & Per-sample/channel standardization after dB transform & (319, 2, 300, 47) \\
vib\_logspec & Pass-mean vibration dB cache & None & Per-sample z-score of dB spectrogram & Per-sample/channel standardization after dB transform & (319, 3, 257, 13) \\
ae\_mel & Pass-mean AE dB cache, inverse-converted to power & 64-bin mel filter bank & Power -> mel sum -> dB & None in cache; no per-sample z-score & (319, 2, 64, 47) \\
vib\_mel & Pass-mean vibration dB cache, inverse-converted to power & 64-bin mel filter bank & Power -> mel sum -> dB & None in cache; no per-sample z-score & (319, 3, 64, 13) \\
\bottomrule
\end{tabular}%
}
\end{table}

\begin{table}[htbp]
\centering
\caption{Compact cache checksum index. Prefixes are the first 12 hexadecimal characters of SHA-256 digests computed from contiguous cached-array bytes. Complete 64-character digests are retained in \path{reports/evidence/tables/cache_provenance.csv} and the public artifact manifest.}
\label{tab:supp-cache-provenance}
\small
\begin{tabular}{@{}lll@{}}
\toprule
Cache & Shape & SHA-256 prefix \\
\midrule
\texttt{ae\_spec} & \texttt{(319, 2, 300, 47)} & \texttt{e19ab882af89} \\
\texttt{vib\_spec} & \texttt{(319, 3, 257, 13)} & \texttt{7f9573cb3fe9} \\
\texttt{ae\_logspec} & \texttt{(319, 2, 300, 47)} & \texttt{33e9bade3f4b} \\
\texttt{vib\_logspec} & \texttt{(319, 3, 257, 13)} & \texttt{eaf5bd9e879f} \\
\texttt{ae\_mel} & \texttt{(319, 2, 64, 47)} & \texttt{ddc30dfd872a} \\
\texttt{vib\_mel} & \texttt{(319, 3, 64, 13)} & \texttt{b1f04f5e4e95} \\
\bottomrule
\end{tabular}
\end{table}


\begin{table}[htbp]
\centering
\caption{Available provenance for the two AE cache planes. The archived
raw-file metadata identify physical acquisition columns, whereas analogue
bandpass settings, gain, sensor model, mounting, and transfer-function
measurements were not retained. Both planes are used in every AE-dB and
AE-dB-z configuration.}
\label{tab:supp-ae-plane-provenance}
\small
\begin{tabular}{@{}cclp{2.5cm}p{4.2cm}@{}}
\toprule
Plane & Raw AE column & Archived label & Cache transformation & Unavailable provenance \\
\midrule
1 & 1 & Narrowband & STFT magnitude $\rightarrow$ dB $\rightarrow$ pass mean & Filter band/order, gain, sensor and mounting response \\
2 & 2 & Broadband & STFT magnitude $\rightarrow$ dB $\rightarrow$ pass mean & Filter band/order, gain, sensor and mounting response \\
\bottomrule
\end{tabular}
\end{table}

\subsection*{Condition-level diagnostics}

Table~\ref{tab:supp-condition-diagnostics} reports per-condition mean
$R_a$ (from the single averaged value per pass), the within-condition
standard deviation of $R_a$, and the mean AE broadband RMS and tri-axial
vibration RMS computed from the cached time-domain features. Condition~7
has the highest mean $R_a$ by a wide margin, but its mean AE RMS and
vibration RMS are not the highest in the design. This pattern is
consistent with condition-specific distribution shift, although it does
not uniquely identify the mechanism.

\begin{table}[htbp]
\centering
\caption{Per-condition mean $R_a$ and mean sensor RMS values.}
\label{tab:supp-condition-diagnostics}
\begin{tabular}{@{}lcccc@{}}
\toprule
Condition & Mean $R_a$ (\(\upmu\)m) & Std $R_a$ (\(\upmu\)m) & AE broad RMS & Vib RMS \\
\midrule
1 & 0.103 & 0.006 & 0.494 & 1.035 \\
2 & 0.089 & 0.004 & 0.579 & 1.531 \\
3 & 0.091 & 0.010 & 0.625 & 1.630 \\
4 & 0.096 & 0.011 & 0.833 & 1.876 \\
5 & 0.097 & 0.005 & 0.743 & 1.458 \\
6 & 0.090 & 0.007 & 0.937 & 1.794 \\
7 & 0.354 & 0.013 & 0.437 & 1.591 \\
8 & 0.242 & 0.011 & 0.301 & 0.896 \\
9 & 0.085 & 0.009 & 0.362 & 1.475 \\
10 & 0.078 & 0.008 & 0.512 & 1.522 \\
11 & 0.108 & 0.005 & 0.306 & 0.747 \\
12 & 0.093 & 0.008 & 0.434 & 1.149 \\
13 & 0.089 & 0.008 & 0.713 & 1.496 \\
14 & 0.106 & 0.005 & 0.793 & 0.789 \\
15 & 0.102 & 0.006 & 0.965 & 1.339 \\
16 & 0.101 & 0.007 & 0.955 & 1.544 \\
\bottomrule
\end{tabular}
\end{table}

\subsection*{LOGO fold target-support analysis}

Table~\ref{tab:supp-logo-extrapolation} shows, for each of the 16 LOGO
folds, the training-set target range and the test-set target range, and
classifies each fold as requiring interpolation or extrapolation.
Because Condition~7 contains by far the highest observed pass-level
$R_a$ values (0.336--0.381~\(\upmu\)m), the training set for all other folds
spans this range, placing all remaining test conditions except
Condition~10 inside the training-label support. Condition~10 is a minor
lower-bound extrapolation because its test minimum is slightly below the
training minimum, but it does not dominate the error budget. The reverse
is true when Condition~7 is held out: the training set maximum drops to
0.259~\(\upmu\)m (Condition~8, the next-roughest condition), so
the evaluated constant-leaf random forest cannot represent targets beyond
the range represented by its terminal-leaf training responses and must
underpredict test targets up to
0.381~\(\upmu\)m by clamping, explaining the large prediction
errors on this fold.

\begin{table}[htbp]
\centering
\caption{LOGO fold target-support analysis. Condition~7 is the only
high-roughness extrapolation fold; Condition~10 is a minor lower-bound
extrapolation.}
\label{tab:supp-logo-extrapolation}
\resizebox{\linewidth}{!}{%
\begin{tabular}{cccccc}
\toprule
Test Fold (Cond. ID) & Train Min $R_a$ (\(\upmu\)m) & Train Max $R_a$ (\(\upmu\)m) & Test Min $R_a$ (\(\upmu\)m) & Test Max $R_a$ (\(\upmu\)m) & Validation Regime \\
\midrule
1 & 0.0660 & 0.3806 & 0.0920 & 0.1128 & Interpolation \\
2 & 0.0660 & 0.3806 & 0.0804 & 0.0994 & Interpolation \\
3 & 0.0660 & 0.3806 & 0.0748 & 0.1052 & Interpolation \\
4 & 0.0660 & 0.3806 & 0.0812 & 0.1190 & Interpolation \\
5 & 0.0660 & 0.3806 & 0.0900 & 0.1086 & Interpolation \\
6 & 0.0660 & 0.3806 & 0.0782 & 0.1050 & Interpolation \\
7 & 0.0660 & 0.2594 & 0.3356 & 0.3806 & Extrap. (High) \\
8 & 0.0660 & 0.3806 & 0.2232 & 0.2594 & Interpolation \\
9 & 0.0660 & 0.3806 & 0.0726 & 0.0998 & Interpolation \\
10 & 0.0726 & 0.3806 & 0.0660 & 0.0978 & Extrap. (Low) \\
11 & 0.0660 & 0.3806 & 0.1002 & 0.1220 & Interpolation \\
12 & 0.0660 & 0.3806 & 0.0796 & 0.1066 & Interpolation \\
13 & 0.0660 & 0.3806 & 0.0782 & 0.1062 & Interpolation \\
14 & 0.0660 & 0.3806 & 0.0990 & 0.1212 & Interpolation \\
15 & 0.0660 & 0.3806 & 0.0940 & 0.1162 & Interpolation \\
16 & 0.0660 & 0.3806 & 0.0916 & 0.1154 & Interpolation \\
\bottomrule
\end{tabular}

}
\end{table}

\subsection*{Input-space novelty and Condition~7/Condition~10 cache contrast}

Table~\ref{tab:supp-input-ood} evaluates a target-free input-space novelty
score retrospectively. For each LOGO fold, the score is the RMS distance
between the held-out condition centroid and its nearest remaining condition
centroid after feature-wise standardisation fitted on the training samples.
The score is associated with the canonical held-out RF error, but the
association was assessed after observing the errors and involves only 16
conditions. It is therefore a candidate input-space monitoring feature, not
a validated or calibrated OOD detector.

\begin{table}[htbp]
\centering
\caption{Condition-level association between a target-free dB-z
input-space novelty score and canonical RF LOGO MAE. The leave-one-condition-
out range recomputes Spearman \(\rho\) after removing each one of the 16
conditions. P-values are unadjusted descriptive values.}
\label{tab:supp-input-ood}
\begin{tabular}{@{}lcc@{}}
\toprule
Score & Spearman \(\rho\) with LOGO MAE & Unadjusted \(p\) \\
\midrule
Nearest dB-z condition-centroid distance & 0.512 & 0.043 \\
Nearest process-parameter distance & -0.256 & 0.339 \\
\bottomrule
\end{tabular}

\vspace{0.4em}
\small Leave-one-condition-out range for the dB-z association: 0.414--0.625.
\end{table}

Table~\ref{tab:supp-condition7-condition10-bands} quantifies the archived
dB-cache contrast between Conditions~7 and~10, which share workpiece speed
and depth of cut but differ in wheel speed. The intervals bootstrap passes
within each condition (20 per condition) and quantify cache-level
differences only; they do not calibrate the sensor response or establish a
causal grinding mechanism.

\begin{table}[htbp]
\centering
\caption{Condition~7 minus Condition~10 archived dB-cache band contrast.
Positive values indicate higher mean dB level for Condition~7.}
\label{tab:supp-condition7-condition10-bands}
\begin{tabular}{@{}lrr@{}}
\toprule
Band & Difference (dB) & Pass-bootstrap 95\% CI (dB) \\
\midrule
AE 0.1--1.0 MHz & 2.77 & [1.37, 4.16] \\
AE $>$1.0 MHz & 2.41 & [1.07, 3.74] \\
Vibration 0.5--2 kHz & -2.94 & [-3.22, -2.67] \\
Vibration 2--15 kHz & -3.14 & [-4.24, -2.09] \\
Vibration $>$15 kHz & -2.33 & [-3.19, -1.48] \\
\bottomrule
\end{tabular}
\end{table}

\subsection*{Out-of-fold SHAP attribution for the current benchmark random forest}

Main Figure~\ref{fig:rf-oof-shap} shows out-of-fold TreeSHAP frequency profiles
under three aggregation rules: sample-pooled absolute attribution,
condition-balanced attribution after normalising each fold to 100\%, and the
same condition-balanced calculation excluding Condition~7 as a descriptive
sensitivity. The latter two prevent folds with larger output deviations or
attribution magnitudes from dominating solely through scale. The sample-pooled
profile remains the direct summary over held-out passes; agreement or movement
across the balanced profiles quantifies whether the high-frequency pattern is
representative across conditions. In the vibration profile, the
condition-balanced mass above 15~kHz is 85.76\%, and it remains 85.74\% after
excluding Condition~7; both peak at 16.2~kHz. The observed upper-band pattern
is therefore not created solely by the difficult fold, although the
sensor-chain caveat near Nyquist remains essential.

The sample-pooled profile reports out-of-fold TreeSHAP mean absolute
importance for the reproducible current random forest trained on fused AE and
vibration dB-z descriptors. A separate RF is fitted in each
LOGO fold after excluding its test and validation conditions, and SHAP is
computed only for that fold's held-out test condition. The 319 held-out
profiles are then pooled by sample count. AE importance is concentrated
above approximately 300~kHz (strongest pooled bin 334~kHz), and vibration
importance is concentrated mainly between approximately 16 and 24~kHz,
with the largest pooled bin at 16.2~kHz. The
high-frequency AE peak is lower than the 1.22~MHz peak reported for the
LightGBM model on dB spectrograms; this shift is consistent with the
dB-z scaling and the random-forest split structure, and
reinforces the need for sensor-transfer-function verification before
treating any single bin as a validated physical signature.

The main figure reports the current reproducible RF and does not attribute the
historical canonical artifact. The global LightGBM and ResNetVibCNN
explanations remain secondary cross-model comparisons.

\subsection*{Seed-repeat results}

Table~\ref{tab:supp-seed-repeat-all} reports the seed-averaged
per-condition LOGO MAE for the equal-budget current random forest and the six
deep-learning baselines that were repeated over five random seeds
(42--46) under the current fixed-setting protocol. These neural runs do not
exactly reproduce the historical tuned checkpoints; they quantify stochastic
variability within the current setting rather than around the historical
canonical models. Condition~7 remains the largest outlier for every model.
The equal-budget random-forest seed-averaged overall MAE is
\(0.0203 \pm 0.0360~\upmu\)m; the DL baselines range from
\(0.0384 \pm 0.0452~\upmu\)m (ResNetVibCNN) to
\(0.0450 \pm 0.0650~\upmu\)m (TrajectoryCNN).

\begin{table}[htbp]
\centering
\caption{Seed-averaged per-condition LOGO MAE ($\upmu$m) for the equal-budget 14-condition random forest and six repeated deep-learning baselines. Headers abbreviate ResNetVibCNN (Vib-CNN), ResNetAECNN (AE-CNN), ChannelAttentionCNN (CA-CNN), ResNetFusion (Fusion), BilinearFusionNetwork (Bilinear), and TrajectoryCNN (Trajectory).}
\label{tab:supp-seed-repeat-all}
\small
\setlength{\tabcolsep}{3.5pt}
\begin{tabular}{@{}cccccccc@{}}
\toprule
Condition & RF & Vib-CNN & AE-CNN & CA-CNN & Fusion & Bilinear & Trajectory \\ \midrule
1 & 0.0054 & 0.0088 & 0.0161 & 0.0104 & 0.0131 & 0.0157 & 0.0080 \\
2 & 0.0032 & 0.0137 & 0.0163 & 0.0115 & 0.0311 & 0.0256 & 0.0634 \\
3 & 0.0094 & 0.0358 & 0.0140 & 0.0139 & 0.0121 & 0.0140 & 0.0110 \\
4 & 0.0083 & 0.0427 & 0.0085 & 0.0165 & 0.0453 & 0.0385 & 0.0119 \\
5 & 0.0036 & 0.0109 & 0.0178 & 0.0091 & 0.0108 & 0.0139 & 0.0046 \\
6 & 0.0061 & 0.0300 & 0.0073 & 0.0110 & 0.0272 & 0.0193 & 0.0182 \\
7 & 0.1509 & 0.1833 & 0.2531 & 0.2614 & 0.2259 & 0.2527 & 0.2567 \\
8 & 0.0585 & 0.1125 & 0.1281 & 0.1006 & 0.1354 & 0.1263 & 0.1434 \\
9 & 0.0092 & 0.0575 & 0.0607 & 0.0524 & 0.0217 & 0.0434 & 0.0567 \\
10 & 0.0105 & 0.0147 & 0.0321 & 0.0308 & 0.0111 & 0.0294 & 0.0387 \\
11 & 0.0082 & 0.0108 & 0.0186 & 0.0164 & 0.0185 & 0.0306 & 0.0123 \\
12 & 0.0100 & 0.0188 & 0.0455 & 0.0463 & 0.0147 & 0.0264 & 0.0069 \\
13 & 0.0123 & 0.0103 & 0.0121 & 0.0210 & 0.0149 & 0.0118 & 0.0125 \\
14 & 0.0142 & 0.0287 & 0.0173 & 0.0165 & 0.0324 & 0.0060 & 0.0612 \\
15 & 0.0099 & 0.0199 & 0.0203 & 0.0235 & 0.0324 & 0.0455 & 0.0067 \\
16 & 0.0050 & 0.0157 & 0.0089 & 0.0359 & 0.0204 & 0.0169 & 0.0081 \\
\bottomrule
\end{tabular}
\end{table}


Paired Wilcoxon signed-rank tests were performed on the 16 condition-level seed-averaged MAEs to describe whether the 14-condition random forest's performance is lower than the repeated 14-condition deep-learning baselines. Table~\ref{tab:supp-seed-repeat-wilcoxon} reports the test statistic, raw p-value, Holm--Bonferroni corrected p-value, and mean MAEs for each comparison. Because the repeated set was selected from the observed benchmark, these results support robustness of the observed RF advantage within the selected repeated-baseline set; they are not independent confirmatory inference.

\begin{table}[htbp]
\centering
\caption{Paired Wilcoxon signed-rank tests comparing the equal-budget
14-condition RF with the six repeated deep-learning baselines. Holm correction
is applied only across this selected six-comparison robustness family.}
\label{tab:supp-seed-repeat-wilcoxon}
\resizebox{\linewidth}{!}{\begin{tabular}{@{}llccccc@{}}
\toprule
Repeated baseline & Input & RF MAE & Baseline MAE & $W$ & $p$ raw & $p$ Holm \\
\midrule
BilinearFusion & \texttt{AE-dB + Vib-dB + phys + PP} & 0.0203 & 0.0447 & 4.0 & 2.136230e-04 & 4.272461e-04 \\
ChannelAttentionCNN & \texttt{AE-dB} & 0.0203 & 0.0423 & 0.0 & 3.051758e-05 & 1.831055e-04 \\
ResNetAECNN & \texttt{AE-dB} & 0.0203 & 0.0423 & 2.0 & 9.155273e-05 & 2.746582e-04 \\
ResNetFusion & \texttt{AE-dB + Vib-dB + phys} & 0.0203 & 0.0417 & 0.0 & 3.051758e-05 & 1.831055e-04 \\
ResNetVibCNN & \texttt{Vib-dB} & 0.0203 & 0.0384 & 1.0 & 6.103516e-05 & 2.441406e-04 \\
TrajectoryCNN & \texttt{Vib trajectory} & 0.0203 & 0.0450 & 13.0 & 2.685547e-03 & 2.685547e-03 \\
\bottomrule
\end{tabular}
}
\end{table}

\begin{table}[htbp]
\centering
\caption{Corrected grouped-validation shallow-MLP sensitivity. The MLP uses
the designated complete validation condition, train-only input and target
standardisation, and a three-candidate training-only hyperparameter grid. The
paired \(p\)-value is descriptive and is not inserted into the historical
14-comparison family.}\label{tab:supp-grouped-mlp}
\begin{tabular}{@{}lrrrr@{}}
\toprule
Model & Mean MAE & Fold SD & Median MAE & Maximum MAE \\
\midrule
Current RF, dB-z & 0.0210 & 0.0361 & 0.0093 & 0.1478 \\
Grouped-validation shallow MLP & 0.0529 & 0.0627 & 0.0275 & 0.2645 \\
\bottomrule
\end{tabular}
\end{table}

\begin{table}[htbp]
\centering
\small
\caption{Corrected grouped-validation shallow-MLP candidate grid. All candidates use training-only input and target standardisation, batch size $\min(32,n_{\mathrm{train}})$, at most 500 epochs, patience 30, and seed 42.}
\label{tab:supp-grouped-mlp-grid}
\begin{tabular}{@{}ccccc@{}}
\toprule
Candidate & Hidden units & Learning rate & L2 $\alpha$ & Activation/optimizer \\
\midrule
1 & 64 & 1e-03 & 1e-03 & ReLU/Adam \\
2 & 128--64 & 1e-03 & 1e-03 & ReLU/Adam \\
3 & 128--64 & 5e-04 & 1e-02 & ReLU/Adam \\
\bottomrule
\end{tabular}
\end{table}

\begin{table}[htbp]
\centering
\small
\caption{Per-fold selection for the corrected grouped-validation shallow MLP. All 48 candidate fits completed without an execution failure; this statement does not imply that every optimiser reached a stationary point.}
\label{tab:supp-grouped-mlp-selections}
\begin{tabular}{@{}ccccc@{}}
\toprule
Test condition & Validation condition & Candidate & Selected epoch & Test MAE ($\upmu$m) \\
\midrule
1 & 8 & 2 & 18 & 0.0129 \\
2 & 5 & 2 & 3 & 0.0305 \\
3 & 14 & 2 & 17 & 0.0397 \\
4 & 16 & 1 & 2 & 0.0278 \\
5 & 12 & 2 & 5 & 0.0156 \\
6 & 9 & 2 & 51 & 0.0113 \\
7 & 14 & 3 & 5 & 0.2645 \\
8 & 5 & 3 & 3 & 0.1376 \\
9 & 7 & 2 & 124 & 0.0185 \\
10 & 11 & 3 & 1 & 0.0712 \\
11 & 3 & 2 & 4 & 0.0699 \\
12 & 7 & 2 & 123 & 0.0239 \\
13 & 11 & 3 & 1 & 0.0271 \\
14 & 11 & 3 & 42 & 0.0254 \\
15 & 8 & 1 & 2 & 0.0264 \\
16 & 5 & 3 & 5 & 0.0445 \\
\bottomrule
\end{tabular}
\end{table}


Table~\ref{tab:supp-seed-repeat-rf} gives the same information for the
random forest only, for readers who want the finest numerical detail.

\begin{table}[htbp]
\centering
\caption{Per-condition seed-averaged LOGO MAE for the equal-budget current random forest (AE-dB-z + Vib-dB-z). Each outer fold uses the same 14 canonical training conditions as the repeated neural baselines; values are averaged over five random seeds.}
\label{tab:supp-seed-repeat-rf}
\begin{tabular}{@{}cc@{}}
\toprule
Condition & Seed-averaged MAE ($\upmu$m) \\
\midrule
1 & 0.0054 \\
2 & 0.0032 \\
3 & 0.0094 \\
4 & 0.0083 \\
5 & 0.0036 \\
6 & 0.0061 \\
7 & 0.1509 \\
8 & 0.0585 \\
9 & 0.0092 \\
10 & 0.0105 \\
11 & 0.0082 \\
12 & 0.0100 \\
13 & 0.0123 \\
14 & 0.0142 \\
15 & 0.0099 \\
16 & 0.0050 \\
\bottomrule
\end{tabular}
\end{table}


\subsection*{Seed-level overall MAEs}

Table~\ref{tab:supp-seed-level} lists the overall LOGO MAE for each of
the five repeated seeds across the equal-budget current random forest and six
deep-learning baselines. The random-forest overall MAE varies by less
than 0.001~\(\upmu\)m across seeds; the deep-learning baselines show
substantially larger seed-to-seed variation.

\begin{table}[htbp]
\centering
\caption{Overall LOGO MAE (\(\upmu\)m) by random seed for the equal-budget
14-condition random forest and six repeated 14-condition deep-learning baselines.}
\label{tab:supp-seed-level}
\begin{tabular}{@{}llccccc@{}}
\toprule
Model & Config & Seed 1 & Seed 2 & Seed 3 & Seed 4 & Seed 5 \\
\midrule
BilinearFusion & \texttt{AE-dB + Vib-dB + physics + PP} & 0.0445 & 0.0419 & 0.0467 & 0.0414 & 0.0492 \\
ChannelAttentionCNN & \texttt{AE-dB} & 0.0503 & 0.0378 & 0.0480 & 0.0334 & 0.0421 \\
RandomForest & \texttt{AE-dB-z + Vib-dB-z} & 0.0210 & 0.0205 & 0.0204 & 0.0201 & 0.0197 \\
ResNetAECNN & \texttt{AE-dB} & 0.0380 & 0.0475 & 0.0478 & 0.0353 & 0.0428 \\
ResNetFusion & \texttt{AE-dB + Vib-dB + physics} & 0.0369 & 0.0380 & 0.0390 & 0.0487 & 0.0459 \\
ResNetVibCNN & \texttt{Vib-dB} & 0.0404 & 0.0283 & 0.0441 & 0.0422 & 0.0369 \\
TrajectoryCNN & \texttt{Vib trajectory} & 0.0439 & 0.0445 & 0.0450 & 0.0452 & 0.0484 \\
\bottomrule
\end{tabular}

\end{table}

\subsection*{Fold-wise attribution stability}

Figure~\ref{fig:supp-shap-stability} shows the dominant vibration and AE
frequency bins from out-of-fold random-forest feature-importance
estimates. The dominant vibration bin is stable around 22~kHz for most
held-out conditions, with a few folds shifting to 17 or 21~kHz. The
dominant AE bin is more variable, shifting between approximately
330~kHz, 880~kHz, and 1.0~MHz depending on the held-out condition.
Because these estimates were obtained with a smaller, faster tree
ensemble for computational reasons, they should be read as a stability
illustration rather than as a definitive per-fold SHAP ranking.

\begin{figure}[htbp]
\centering
\includegraphics[width=\linewidth]{shap_fold_stability.png}
\caption{Dominant frequency bins from out-of-fold random-forest
impurity-based feature-importance estimates. \textbf{a}, Vibration.
\textbf{b}, AE.}\label{fig:supp-shap-stability}
\end{figure}

\subsection*{RF OOB coverage by condition}

Table~\ref{tab:rf-oob-coverage} reports the empirical coverage of the
95\% OOB-calibrated intervals for a separate 15-condition RF uncertainty
variant, broken down by grinding condition. Coverage is highly uneven:
Condition~7 has 0\% coverage, while several conditions reach 100\%.
This pattern motivates the condition-level calibration discussion in the
main text.

\begin{table}[htbp]
\centering
\caption{Per-condition descriptive target coverage of nominal 95\%
OOB-calibrated intervals for a separate 15-condition RF uncertainty variant
(AE-dB-z + Vib-dB-z).}
\label{tab:rf-oob-coverage}
\begin{tabular}{@{}lcccc@{}}
\toprule
Condition & Passes & Coverage (\%) & Mean width ($\upmu$m) & MAE ($\upmu$m) \\
\midrule
1 & 19 & 100.0\% & 0.0288 & 0.0054 \\
2 & 20 & 100.0\% & 0.0338 & 0.0031 \\
3 & 20 & 65.0\% & 0.0307 & 0.0103 \\
4 & 20 & 95.0\% & 0.0776 & 0.0081 \\
5 & 20 & 100.0\% & 0.0277 & 0.0030 \\
6 & 20 & 80.0\% & 0.0198 & 0.0060 \\
7 & 20 & 0.0\% & 0.1973 & 0.1368 \\
8 & 20 & 100.0\% & 0.2909 & 0.0681 \\
9 & 20 & 100.0\% & 0.0878 & 0.0108 \\
10 & 20 & 85.0\% & 0.0592 & 0.0105 \\
11 & 20 & 90.0\% & 0.0394 & 0.0090 \\
12 & 20 & 75.0\% & 0.0247 & 0.0090 \\
13 & 20 & 40.0\% & 0.0240 & 0.0131 \\
14 & 20 & 90.0\% & 0.1284 & 0.0142 \\
15 & 20 & 65.0\% & 0.0278 & 0.0096 \\
16 & 20 & 100.0\% & 0.0284 & 0.0060 \\
\bottomrule
\end{tabular}

\end{table}

\subsection*{Measurement uncertainty caveat}

Only the manufacturer-quoted profilometer repeatability (0.01~\(\upmu\)m) is
available. The raw profilometer trace files needed to compute an
experimental repeatability or gauge R\&R were not archived. The best
reported MAE (0.020~\(\upmu\)m) is therefore close to the manufacturer
repeatability, so differences smaller than roughly 0.005--0.010~\(\upmu\)m
should be interpreted cautiously.

\subsection*{Deep Learning Baselines Architecture Details}

To support replication, Table~\ref{tab:supp-dl-arch} details the architectural configurations, layer dimensions, activation functions, regularization (dropout), and parameter counts for the nine architecturally documented deep-learning baselines. The 38-input canonical LOGO inventory contains 25 exactly reconstructable model variants plus two archival variants with incomplete resolved-configuration provenance. The latter are retained only as historical inventory entries and are not central to any conclusion.

\begin{table}[htbp]
\centering
\caption{Detailed architectural specifications and hyperparameter configurations for the nine evaluated deep-learning baselines.}
\label{tab:supp-dl-arch}
\begin{tabular}{lp{4cm}ccc}
\toprule
Model Baseline & Key Layer Blocks & Hidden Dim / Embed Dim & Dropout $p$ & Parameter Count \\
\midrule
\texttt{ResNetAECNN} & 2D Conv stem, 2 SE-ResNet blocks, GAP, MLP regressor & 64 & 0.3 & 86,402 \\
\texttt{ResNetVibCNN} & 2D Conv stem, 2 SE-ResNet blocks, GAP, MLP regressor & 64 & 0.3 & 86,546 \\
\texttt{ResNetFusion} & Dual ResNet encoders, feature concatenation, MLP regressor & 128 & 0.3 & 167,506 \\
\texttt{BilinearFusionNetwork} & Dual ResNet encoders, low-rank bilinear pooling layer, MLP & 64 & 0.3 & 171,425 \\
\texttt{TrajectoryCNN} & 1D sequence CNN stem, 2 temporal blocks, GAP, MLP regressor & 64 & 0.5 & 47,905 \\
\texttt{GraphNeuralNetworkFusion} & 4 Modality MLP encoders, 2 Graph Attention (GAT) layers, GAP, MLP & 32 & 0.3 & 152,673 \\
\texttt{TabTransformerRegressor} & Feature-wise linear embedding, 2 Transformer encoder layers, MLP & 16 & 0.3 & 6,353 \\
\texttt{ChannelAttentionCNN} & 2D Conv, CBAM block (Channel + Spatial attention), GAP, MLP & 64 & 0.3 & 30,216 \\
\texttt{MultiscaleSpectrogramCNN} & Multi-resolution 2D Conv filter banks, GAP, MLP regressor & 128 & 0.3 & 34,466 \\
\bottomrule
\end{tabular}
\end{table}


\subsection*{Complete model and input inventory}
\begin{longtable}{@{}p{0.42\linewidth}cp{0.43\linewidth}@{}}
\caption{Complete LOGO-only model inventory. The archived canonical benchmark contains 27 model variants evaluated across 38 distinct input configurations.}\label{tab:supp-model-inventory}\\
\toprule
Model variant & Inputs ($n$) & Configuration IDs \\
\midrule
\endfirsthead
\toprule
Model variant & Inputs ($n$) & Configuration IDs \\
\midrule
\endhead
\texttt{AttentionMLP} & 2 & \texttt{C08, C09} \\
\texttt{BilinearFusionNetwork} & 19 & \texttt{C03, C04, C06, C07, C08, C09, C16, C17, C18, C19, C20, C21, C25, C26, C27, C30, C31, C35, C36} \\
\texttt{ChannelAttentionCNN} & 1 & \texttt{C16} \\
\texttt{CrossAttentionFusion} & 1 & \texttt{C08} \\
\texttt{FeatureOnlyModel} & 2 & \texttt{C08, C09} \\
\texttt{GraphNeuralNetworkFusion} & 19 & \texttt{C03, C04, C06, C07, C08, C09, C16, C17, C18, C19, C20, C21, C25, C26, C27, C30, C31, C35, C36} \\
\texttt{HierarchicalAttentionNetwork} & 1 & \texttt{C09} \\
\texttt{LightGBMModel} & 20 & \texttt{C03, C05, C06, C08, C09, C10, C11, C12, C13, C14, C15, C16, C18, C21, C27, C30, C32, C33, C34, C35} \\
\texttt{MeanTrajectoryMLP} & 5 & \texttt{C22, C23, C24, C37, C38} \\
\texttt{MultiHeadAttentionFusion} & 2 & \texttt{C08, C09} \\
\texttt{MultiscaleSpectrogramCNN} & 1 & \texttt{C35} \\
\texttt{ParamsMLP} & 1 & \texttt{C27} \\
\texttt{PhysicsGuidedAttentionNetwork} & 1 & \texttt{C09} \\
\texttt{PhysicsInformedFusionNet} & 1 & \texttt{C25} \\
\texttt{PhysicsMLP} & 8 & \texttt{C03, C06, C08, C09, C18, C21, C27, C30} \\
\texttt{RandomForestModel} & 20 & \texttt{C03, C05, C06, C08, C09, C10, C11, C12, C13, C14, C15, C16, C18, C21, C27, C30, C32, C33, C34, C35} \\
\texttt{ResNetAECNN} & 1 & \texttt{C16} \\
\texttt{ResNetFusion} & 9 & \texttt{C16, C17, C18, C19, C20, C21, C25, C35, C36} \\
\texttt{ResNetVibCNN} & 1 & \texttt{C35} \\
\texttt{RidgeRegressionModel} & 5 & \texttt{C09, C16, C18, C27, C35} \\
\texttt{SSLResNetVibCNN} & 1 & \texttt{C35} \\
\texttt{ShallowMLPModel} & 3 & \texttt{C09, C27, C35} \\
\texttt{TabTransformerRegressor} & 12 & \texttt{C03, C04, C05, C06, C07, C08, C09, C26, C27, C30, C31, C32} \\
\texttt{TrajectoryAttention} & 5 & \texttt{C22, C23, C24, C37, C38} \\
\texttt{TrajectoryCNN} & 5 & \texttt{C22, C23, C24, C37, C38} \\
\texttt{TransferFeatureMLP} & 4 & \texttt{C01, C02, C28, C29} \\
\texttt{XGBoostModel} & 12 & \texttt{C03, C05, C06, C08, C09, C16, C18, C21, C27, C30, C32, C35} \\
\bottomrule
\end{longtable}

\begin{longtable}{@{}cp{0.82\linewidth}@{}}
\caption{Configuration key for Table~\ref{tab:supp-model-inventory}.}\\
\toprule
ID & Input configuration \\
\midrule
\endfirsthead
\toprule
ID & Input configuration \\
\midrule
\endhead
C01 & AE embedding + Vib embedding \\
C02 & AE embedding + Vib embedding + PP \\
C03 & AE features \\
C04 & AE features + physics features \\
C05 & AE features + PP \\
C06 & AE features + Vib features \\
C07 & AE features + Vib features + physics features \\
C08 & AE features + Vib features + physics features + PP \\
C09 & AE features + Vib features + PP \\
C10 & AE-dB-z \\
C11 & AE-dB-z + Vib-dB-z \\
C12 & AE-dB-z + Vib-dB-z + PP \\
C13 & AE-log-mel \\
C14 & AE-log-mel + Vib-log-mel \\
C15 & AE-log-mel + Vib-log-mel + PP \\
C16 & AE-dB \\
C17 & AE-dB + physics features \\
C18 & AE-dB + Vib-dB \\
C19 & AE-dB + Vib-dB + physics features \\
C20 & AE-dB + Vib-dB + physics features + PP \\
C21 & AE-dB + Vib-dB + PP \\
C22 & AE trajectory \\
C23 & AE trajectory + Vib trajectory \\
C24 & AE trajectory + Vib trajectory + PP \\
C25 & all inputs \\
C26 & physics features + PP \\
C27 & PP \\
C28 & Vib embedding \\
C29 & Vib embedding + PP \\
C30 & Vib features \\
C31 & Vib features + physics features \\
C32 & Vib features + PP \\
C33 & Vib-dB-z \\
C34 & Vib-log-mel \\
C35 & Vib-dB \\
C36 & Vib-dB + physics features \\
C37 & Vib trajectory \\
C38 & Vib trajectory + PP \\
\bottomrule
\end{longtable}


Table~\ref{tab:supp-model-constructors} is a print-readable summary of the
resolved model inventory. Complete signatures and pair-level configuration
notes remain in the hashed machine-readable JSON records cited in its caption.
Neural models use the common optimisation protocol stated in
Section~\ref{predictive-models}; input IDs determine which constructor branches
are active.

\begingroup
\footnotesize
\setlength{\tabcolsep}{3.5pt}
\begin{longtable}{@{}p{0.23\linewidth}p{0.20\linewidth}p{0.39\linewidth}p{0.12\linewidth}@{}}
\caption{Concise resolved inventory of 25 reconstructable variants and two archival records. Input IDs are defined in Table~\ref{tab:supp-model-inventory}. Complete signatures and pair-level notes are in \texttt{canonical\_model\_constructor\_defaults.json} and \texttt{resolved\_model\_configuration\_pairs.json} (SHA-256 prefix \texttt{aff13bdd271d}) at public artifact commit \texttt{744631057151}. Complete identifiers are retained in the public manifest. The two archival entries require Vib-dB channel or modality overrides that were not retained. The historical shallow MLP used an internal random 10\% validation split; the corrected grouped-validation sensitivity is reported separately.}
\label{tab:supp-model-constructors}\\
\toprule
Model & Configured input & Key exposed defaults & Status \\
\midrule
\endfirsthead
\toprule
Model & Configured input & Key exposed defaults & Status \\
\midrule
\endhead
\texttt{AttentionMLP} & 2 configurations & dropout=0.3 & Reconstructable \\
\texttt{BilinearFusionNetwork} & 19 configurations & embed\_dim=32, rank=16, dropout=0.3 & Reconstructable \\
\texttt{ChannelAttentionCNN} & AE-dB & in\_channels=2, dropout=0.3 & Reconstructable \\
\texttt{CrossAttentionFusion} & AE features + Vib features + physics features + PP & embed\_dim=32, num\_heads=2, num\_layers=2, dropout=0.3 & Reconstructable \\
\texttt{FeatureOnlyModel} & 2 configurations & No exposed architecture arguments & Reconstructable \\
\texttt{GraphNeuralNetworkFusion} & 19 configurations & embed\_dim=32, hidden\_dim=32, num\_heads=2, num\_layers=2 & Reconstructable \\
\texttt{HierarchicalAttentionNetwork} & AE features + Vib features + PP & embed\_dim=32, num\_heads=2, num\_layers=1, dropout=0.3 & Reconstructable \\
\texttt{LightGBMModel} & 20 configurations & n\_estimators=200, max\_depth=5, learning\_rate=0.05, subsample=0.8 & Reconstructable \\
\texttt{MeanTrajectoryMLP} & 5 configurations & ae\_input\_dim=6, vib\_input\_dim=9, hidden\_dim=128, dropout=0.3 & Reconstructable \\
\texttt{MultiHeadAttentionFusion} & 2 configurations & embed\_dim=32, num\_heads=2, num\_layers=2, dropout=0.3 & Reconstructable \\
\texttt{MultiscaleSpectrogramCNN} & Vib-dB & in\_channels=2, base\_channels=16, dropout=0.3 & Archival \\
\texttt{ParamsMLP} & PP & No exposed architecture arguments & Reconstructable \\
\texttt{PhysicsGuidedAttentionNetwork} & AE features + Vib features + PP & embed\_dim=32, dropout=0.3 & Reconstructable \\
\texttt{PhysicsInformedFusionNet} & all inputs & embed\_dim=64, num\_heads=4, num\_layers=2, dropout=0.3 & Reconstructable \\
\texttt{PhysicsMLP} & 8 configurations & No exposed architecture arguments & Reconstructable \\
\texttt{RandomForestModel} & 20 configurations & n\_estimators=200, max\_depth=8, n\_jobs=1, random\_state=42 & Reconstructable \\
\texttt{ResNetAECNN} & AE-dB & dropout=0.3, hidden\_dim=64 & Reconstructable \\
\texttt{ResNetFusion} & 9 configurations & embed\_dim=64, dropout=0.3 & Reconstructable \\
\texttt{ResNetVibCNN} & Vib-dB & dropout=0.3, hidden\_dim=64 & Reconstructable \\
\texttt{RidgeRegressionModel} & 5 configurations & alpha=1.0, random\_state=42 & Reconstructable \\
\texttt{SSLResNetVibCNN} & Vib-dB & modality='ae', embed\_dim=64, proj\_dim=32, dropout=0.3 & Archival \\
\texttt{ShallowMLPModel} & 3 configurations & hidden\_layer\_sizes=(128, 64), alpha=0.0001, learning\_rate\_init=0.001, max\_iter=500 & Reconstructable \\
\texttt{TabTransformerRegressor} & 12 configurations & max\_features=49, embed\_dim=16, num\_heads=2, num\_layers=2 & Reconstructable \\
\texttt{TrajectoryAttention} & 5 configurations & ae\_input\_dim=6, vib\_input\_dim=9, hidden\_dim=64, dropout=0.3 & Reconstructable \\
\texttt{TrajectoryCNN} & 5 configurations & ae\_input\_dim=6, vib\_input\_dim=9, channels=None, kernel\_size=5 & Reconstructable \\
\texttt{TransferFeatureMLP} & 4 configurations & embed\_dim=512, hidden\_dim=128, dropout=0.3, use\_ae=True & Reconstructable \\
\texttt{XGBoostModel} & 12 configurations & n\_estimators=200, max\_depth=5, learning\_rate=0.05, subsample=0.8 & Reconstructable \\
\bottomrule
\end{longtable}
\endgroup

